\section{Introduction}
\label{sec:intro}

Bid-rigging cartels in public procurement face a practical constraint:
they must simulate competition to avoid detection. A common solution is
\emph{cover bidding}---deploying firms that submit losing bids to
legitimize the process and create the appearance of rivalry
\citep{marshall2012economics, asker2010leniency}. Detecting such
behavior remains hard, especially where bid-level data are unavailable
\citep{harrington2008detecting, imhof2019detecting}.

This paper proposes a firm-level diagnostic built on a participation
anomaly: firms that repeatedly enter tenders but never win. We call
them \emph{frequent losers} (FL). Sustained participation with a zero
win rate is difficult to reconcile with profit-maximizing behavior
when bidding is costly, and matches the profile of cover bidders
documented in prosecuted cartels---and the screen requires only
participation and outcome records, making it applicable where the
bid-level data demanded by existing tools are incomplete or absent.

Our laboratory is S\~ao Paulo's Bolsa Eletr\^onica de Compras (BEC),
covering 4.5 million tender-items and 40 million bids over
2009--2019. Among 41,000 firms, 16,843 never win a single tender.
An IQR-based outlier rule applied to this always-loser population
identifies 2,735 whose participation frequency cannot be squared with
profit-maximizing entry---the frequent losers. A firm that bids on
50 tenders over a decade without winning a single one is either
extraordinarily unlucky or playing a different game.

Within-item, within-year, within-purchasing-unit comparisons show a
3.6--7.7\% conditional price gap in FL-present tenders. The range
spans temporal cross-fitting (3.6\%), OLS with high-dimensional fixed
effects (6.4\%), and matching (CEM: 7.7\%; IPW: 5.5\%). After
absorbing item fixed effects the raw gap collapses from 2.43 to 0.050
log points (overlap 0.978; Table~\ref{tab:conditional_descstats},
Appendix~\ref{app:competition}), so the regressions compare like with
like. Comparability, however, is not identification.

These are conditional associations, not causal effects. The gap may
reflect collusive environments or unobserved market characteristics
that jointly attract FL firms and higher prices; we cannot fully
separate the two. Sensitivity analysis
(\citealt{cinelli2020making}, $RV = 17.5\%$) bounds the scope for
omitted confounders, and the event-study pre-trends are candidly
presented as an unresolved ambiguity
(Section~\ref{sec:additional_id}). The contribution is primarily
\emph{diagnostic}: FL presence identifies procurement environments
systematically associated with higher prices and with known cartel
activity.

An external consistency check against CADE (Brazil's competition
authority) shows FL firms are 3.5$\times$ more present in
cartel-affected markets ($p < 0.001$). This excess reflects market
proximity rather than direct cartel membership; its nuances are taken
up in Section~\ref{sec:cade}. The price association barely moves when
we drop all CADE-involved markets ($\hat{\beta} = 0.062$), so the
marker picks up something beyond the enforcement record.

Several mechanism diagnostics point toward cover bidding as one
plausible driver without being individually or jointly conclusive.
Working from the broadest pattern to the most specific: the price
effect concentrates in competitive markets, where a cover-bidding
framework predicts such activity is most profitable;
bid-coordination tests reject exchangeability of FL and non-FL
residuals; structural estimation favors coordinated cover bidding
($\hat{\sigma}_c / \hat{\sigma}_g = 0.72$). The sharpest
mechanism-consistent heterogeneity comes from convite tenders where
the minimum-bidder rule does not bind: the FL coefficient there is
7.6\%, larger than in tenders where the rule forces participation,
though selection within the non-binding subsample cannot be fully
excluded (Section~\ref{sec:additional_id}).

\subsection{Contributions}

\emph{First}, a firm-level diagnostic for suspicious procurement
environments. The FL screen operates at the firm level, needs no bid
values, and provides information orthogonal to bid-dispersion screens
(horse-race correlation 0.060). FL flags triage at scale; bid-level
analysis follows on flagged tenders.

\emph{Second}, a robust empirical regularity: the price gap survives
four estimators that differ in functional form but share the
selection-on-observables assumption. The cross-fit decomposition shows
that most of the OLS--cross-fit discrepancy is classification noise,
not mechanical bias---tightening the informative range.

\emph{Third}, evidence consistent with cover bidding as one potential
driver, assembled with explicit acknowledgment that observational data
cannot close the identification. The convite voluntary-deployment
premium, the Bajari--Ye results, and the structural regime
identification narrow the space of non-collusive explanations without
over-claiming causation. From an IO perspective, the finding that
participation patterns alone can flag suspicious environments suggests
that firm-level entry behavior is itself an informative equilibrium
object---one that existing screens, focused on bid distributions,
have overlooked.

After describing the data and FL classification
(Sections~\ref{sec:data}--\ref{sec:structural}), we present the
empirical strategy (Section~\ref{sec:empirical}), the CADE
consistency check (Section~\ref{sec:cade}), the price regressions
and mechanism diagnostics
(Sections~\ref{sec:results}--\ref{sec:mechanisms}), robustness
(Section~\ref{sec:robustness}), and a conclusion
(Section~\ref{sec:conclusion}).
